\section*{Chapter \ref{ch:pl:cloud-against-on-premises} --- Cloud Against On-Premises}

\begin{solution}{exo:pl:cloud-against-on-premises:1}
The instance's hourly price divided by its 96 threads: 4.46 cents on demand, 1.96 cents reserved for three years without upfront payment, and
1.81 cents on spot (on 28 September 2026).
\end{solution}

\begin{solution}{exo:pl:cloud-against-on-premises:2}
Above 1.886/4.284, 44\%: below that, paying every hour for \omterm{def:pl:cloud-against-on-premises:capacity}{reserved capacity} costs more than paying on demand only for the hours used.
\end{solution}

\begin{solution}{exo:pl:cloud-against-on-premises:3}
\$29\,491: 10~TB at 9 cents a gigabyte, 40~TB at 8.5, 100~TB at 7 and the remaining 350~TB at 5 (with 1\,024 GB to the TB). Bringing it in was
free.
\end{solution}

\begin{solution}{exo:pl:cloud-against-on-premises:4}
Both are paid every hour whether used or not, so for any unit of capacity the cheaper hourly price wins at every utilisation: 0.40 cents owned
against 1.96 reserved. \omterm{def:pl:cloud-against-on-premises:capacity}{Reserved capacity} wins only when owning is not possible --- no data centre, no time to buy, a need that lasts less than
the server's life.
\end{solution}

\begin{solution}{exo:pl:cloud-against-on-premises:5}
The sweep's length is set by its longest backtests, 7.5 hours each. An interruption restarts one from the beginning, so if it strikes six hours
in, that backtest ends after 13.5 hours and the sweep with it: 14.4 hours at the advisor's band.
\end{solution}

\begin{solution}{exo:pl:cloud-against-on-premises:6}
The owned server's cost per unit of work halves, and so does the break-even: 4.5\% against on demand, 11\% against spot. Speed matters as much
as price, which is why the comparison should be made in backtests per dollar, measured.
\end{solution}

\begin{solution}{exo:pl:cloud-against-on-premises:7}
The sweep adds 2\,400 threads to every weeknight's batch, a level used 24\% of hours, so the planner owns 7\,900 threads instead of 5\,500 and the
year costs \$291\,003 with on-demand bursts, against \$361\,215: moving flexible work into the hours owned capacity already covers raises its
utilisation and saves more than any price negotiation.
\end{solution}

\begin{solution}{exo:pl:cloud-against-on-premises:8}
\omterm{def:pl:cloud-against-on-premises:capacity}{Spot capacity} can be reclaimed at any moment: interactive work and anything without checkpoints loses its progress, and the chapter's sweep
without checkpoints took twice as long at the advisor's band. And for the base load, used all the time, owned capacity is more than four times cheaper per
thread-hour than spot at the model's prices.
\end{solution}

\begin{solution}{pb:pl:cloud-against-on-premises:1}
\begin{enumerate}
\item 4.46 cents on demand, 2.95 reserved one year, 1.96 reserved three years, 1.81 spot: the provider's published price list of 25 September
2026 and its spot feed of 28 September, divided by 96 threads.
\item Cited: the CPU's list price and the average commercial electricity price. Assumed: the rest of the server (as much again as the CPUs),
four-year life, 1.3~kW, overhead 1.4, \$300 of space and a two-hundredth of an engineer a month.
\item \$1\,508 a month, 0.40 cents per thread-hour fully used.
\item Fixed capacity (owned, reserved) costs its hourly price divided by utilisation; metered capacity (on demand, spot) costs its price per
hour used.
\item 9.0\% against on demand and 22.3\% against spot; 14.5\% and 35.8\% at twice the server price, 20.0\% and 49.2\% at three times.
\item Base 1\,500 threads, 1\,000 more in weekday hours, 4\,000 more on weeknights, 12\,000 more for eight hours on Saturdays, twenty bursts of
8\,000 for four hours; mean 3\,364, peak 13\,500.
\item From the load-duration curve: the $k$-th unit is used in the share $f(k)$ of hours, and it is owned if owning costs less than $f(k)$ times
the hourly price, reserved on the same test when owning is excluded, otherwise bought by the hour.
\item 5\,500 threads owned; 496 hours have demand above it.
\item \$1\,315\,233 on demand only, \$940\,460 reserved and on demand, \$361\,215 owned and on demand, \$281\,883 owned with spot bursts.
\item Because the demand has a base used most of the time and peaks used rarely, on opposite sides of the break-even.
\item Savings of 65\% over on demand over the last 30 days, and an interruption frequency in the highest band, above 20\% in the trailing
month.
\item At the advisor's band, \$94 and 14.4 hours; at 0.2 an hour, \$185 and 59 hours; against \$203 and 7.5 hours on demand.
\item The loss is bounded at one chunk: \$80--82 and under 9 hours at every rate.
\item \$11\,776 a month in standard object storage; \$29\,491 to take a copy out once.
\item In the EU, under the EBA's outsourcing guidelines, among other things a documented exit strategy for a critical or important function.
\item \textbf{Named result.} At the cited prices and the model's owned server, owning costs less than renting on demand above 9\% utilisation
(22\% against spot; 20\% and 49\% even at three times the server price). Chapter~13's sweep costs \$203 on demand and \$82 on spot without
interruptions; at the provider's highest interruption band it costs \$94 and takes 14.4 hours instead of 7.5, and with 15-minute checkpoints
\$80 in under 9 hours.
\item The relative speed of the two machines on the firm's own backtests, measured; then the server's full price from a quote.
\item Where most of the computation that reads it runs; with one authoritative copy, and computation sent to it.
\item Demand dominated by rare large peaks, no data centre or staff, or data already in the cloud.
\item The shape of the demand against the price of fixed and metered capacity, where the data is, and how the work can be interrupted.
\end{enumerate}
\end{solution}

\begin{solution}{iq:pl:cloud-against-on-premises:1}
On demand: no commitment, highest price, for irregular work. Reserved: one or three years committed at a lower price, for steady load. Spot:
spare capacity at the lowest price that can be reclaimed, for interruptible, checkpointed batch work.
\iqlookfor{Commitment, price, interruptibility.}
\end{solution}

\begin{solution}{iq:pl:cloud-against-on-premises:2}
Build the demand profile and its load-duration curve, price owned capacity per hour with its assumptions stated, and buy each level of capacity
the cheapest way given the share of hours it is used; check data location and interruptibility; compare in work per dollar, measured.
\iqlookfor{Utilisation and break-even, not list prices.}
\end{solution}

\begin{solution}{iq:pl:cloud-against-on-premises:3}
Make every task restartable: checkpoint long tasks at intervals, keep tasks idempotent with results written atomically, let the scheduler
resubmit interrupted work, and spread over instance types and zones.
\iqlookfor{Checkpoints and idempotent tasks.}
\end{solution}

\begin{solution}{iq:pl:cloud-against-on-premises:4}
Moving it in is free but it must then be stored there (about \$12\,000 a month in standard storage) and taking it out costs about \$29\,000 a
copy; decide where the authoritative copy lives and send computation to it, or copy only the subset the research needs.
\iqlookfor{\omterm{def:pl:cloud-against-on-premises:gravity}{Data gravity} and data-transfer charges.}
\end{solution}

\begin{solution}{iq:pl:cloud-against-on-premises:5}
Owned cluster sized at the level used above the break-even share of hours; one scheduler over owned and cloud nodes; bursts on spot for
checkpointed work and on demand for the rest; data near the owned cluster with cached subsets in the cloud; everything as code; an exit plan.
\iqlookfor{Size owned for the base, burst the peaks.}
\end{solution}
